\documentclass[11pt]{article}
\usepackage[margin=1.15in]{geometry}
\usepackage[T1]{fontenc}
\usepackage{mathpazo}
\usepackage{microtype}
\usepackage{amsmath,amssymb,amsthm}
\usepackage{booktabs}
\usepackage{array}
\usepackage[round,authoryear]{natbib}
\usepackage{setspace}
\usepackage{enumitem}
\usepackage[colorlinks=true,linkcolor=black,citecolor=black,urlcolor=black]{hyperref}
\setstretch{1.08}
\setlength{\parskip}{0.35em}

\newtheorem{hypothesis}{Hypothesis}
\newtheorem{definition}{Definition}
\theoremstyle{remark}
\newtheorem{remark}{Remark}

\title{Transformational Familiarity:\\Food Neophobia, Autonomy, and the Construction of Taste}
\author{Flyxion\\\small Independent Researcher}
\date{October 2026}

\begin{document}
\maketitle

\begin{abstract}
\noindent Children who refuse unfamiliar foods are usually described as picky, a label that treats a heterogeneous set of behaviors as a single defect to be corrected. This essay argues that the label obscures at least five distinct processes: neophobia toward novelty, sensory aversion to particular properties, a stable preference for familiar foods, ordinary variation in appetite, and resistance to coercion. These produce similar behavior at the table but call for different responses, and some of them are made worse by the interventions that work for others. Drawing on the literature on exposure, social modeling, responsive feeding, and food categorization, the essay develops a framework in which acceptance of a food depends on how that food is connected to the child's existing categories and on how much control the child retains over the pace of approach. It proposes the Transformational Familiarity Hypothesis: under safe, voluntary, and pace-controlled conditions, an unfamiliar food becomes more approachable when it is introduced through a sequence of preparations that preserve recognizable features of an already accepted food. The hypothesis is formalized, distinguished from the established findings on repeated exposure, and paired with a worked example, a typology of sensory, procedural, and causal familiarity, and a protocol by which it could be tested. Its conclusion is that food acceptance depends in part on the flexibility of the surrounding food environment, and not solely on the flexibility of the eater.
\end{abstract}

\tableofcontents
\newpage

\section{Introduction}\label{sec:intro}

Few childhood behaviors generate as much household conflict, and as little agreement about cause, as the refusal of food. Parents, clinicians, and grandparents hold confident and mutually incompatible views about what picky eating is and what should be done about it. Some regard it as a discipline problem to be settled by firmness, others as a temperament to be waited out, others as a sensory difference that must be accommodated, and others as a symptom that the child has been given too many alternatives. Each of these views is supported by cases in which it appears to have worked, and each fails in cases that the others handle easily. A framework that accounts for this pattern of partial success is more useful than another prescription.

The central claim of this essay is that picky eating is best understood as a problem about how foods are introduced, classified, and made available for exploration, and only secondarily as a property of the child. A refusal is a response to a particular food, in a particular form, presented under particular conditions, and each of those three elements can be altered by the adults who organize the meal. This does not deny that children differ in temperament, sensory sensitivity, and the strength of their neophobia. A twin study of children aged 8 to 11 estimated the heritability of food neophobia at 0.78 \citep{cooke2007genetic}. Heritability describes the share of variation in a trait that is attributable to genetic differences within a particular population and environment. It does not mean that 78 percent of any individual child's neophobia is genetically determined, and it does not set a limit on the influence that changes in the environment can have. Some children also have feeding difficulties of clinical significance. The claim is that, within the range of ordinary variation, the structure of the food environment contributes to the outcome, and that this contribution is both larger and more tractable than the vocabulary of defect and discipline suggests.

The contribution of the essay can be located precisely. The repeated-exposure literature is concerned with changing the eater's response to a fixed food, by presenting the same item until it becomes familiar. The proposal developed here is concerned with changing the pathway by which the eater encounters the food, and treats the history of encounters, together with the form of the food and the child's state, as what determines the response. The first is a well-supported account of one mechanism. The second is a conjecture about a different and possibly complementary mechanism, and separating it from the first is the main burden of the formal development and the experimental protocol.

The essay proceeds in four movements. The first (Sections~\ref{sec:disagg} and~\ref{sec:evidence}) separates the phenomena that are grouped under the picky-eating label and reviews what the empirical literature establishes about exposure, modeling, and pressure. The second (Sections~\ref{sec:categories} and~\ref{sec:hypothesis}) develops a theory of food categories and proposes the Transformational Familiarity Hypothesis, a conjecture about the use of intermediate preparations to connect accepted and unaccepted foods. The third (Sections~\ref{sec:approach} through~\ref{sec:economics}) examines the conditions under which such transformations are likely to work: observational and participatory involvement, procedural and causal familiarity gained through cultivation and cooking, a distribution of agency in which adults structure the environment and children retain authority over participation, and the culinary competence that makes a flexible environment affordable. The fourth (Sections~\ref{sec:limits} through~\ref{sec:conclusion}) states the limits of the argument, proposes an experimental protocol, and concludes.

A note on status is necessary at the outset. The established literature on repeated exposure and social modeling is substantial and is reviewed here as such. The central hypothesis of the essay, by contrast, is a proposal, not a finding. Observations from ordinary cooking suggest that foods connected by overlapping properties are easier to approach than foods presented without such connections, but observations of this kind are vulnerable to every familiar confound: the child may have simply grown older, the adult may have simply become more skilled, and repeated exposure may have done the work attributed to the transformation. The purpose of the formal apparatus developed below is to make the proposal precise enough that it can fail.

\section{Picky Eating Is Not a Unitary Phenomenon}\label{sec:disagg}

The terms picky eating, fussy eating, and food fussiness are used in the research literature and in everyday speech with considerable looseness. \citet{dovey2008food} reviewed the field and distinguished food neophobia, the reluctance to eat new foods, from picky or fussy eating, which involves rejecting a substantial range of foods that are familiar as well as unfamiliar, while noting that the two constructs are inter-related and share some influences. Behaviors that look alike at the table therefore need not share a cause, and the labels are applied inconsistently in everyday speech. The first task of any serious account is therefore to disaggregate the category.

Five distinct processes can be separated analytically, with the understanding that they overlap in practice. The first is \emph{food neophobia}, the reluctance to eat or try novel foods, which has been measured in adults with an established scale \citep{pliner1992development} and in children with parent-report and behavioral measures \citep{pliner1994development}. In an omnivore, neophobia is understood as one side of a balance. An animal that eats many kinds of food needs to explore unfamiliar items in order to meet its nutritional needs, and needs caution about them in order to avoid ingesting toxins \citep{rozin1976selection}. Neophobia is accordingly an adaptive disposition, and its strength in early childhood coincides with the period in which children become mobile and independent enough to ingest things their caregivers have not vetted.

The second process is \emph{sensory aversion}, a negative response to a specific property of a food, such as a texture, a temperature, a bitterness, or a smell. This differs from neophobia because it can attach to foods the child has encountered many times, and because it need not diminish with exposure. Food selectivity is common in children with autism spectrum disorders, and a narrative review discusses sensory sensitivity as a possible contributor to it \citep{cermak2010food}. The distinction has practical weight. A child who rejects a novel food because it is novel is likely to respond to repeated, low-pressure contact. A child who rejects a food because its texture is aversive may not, and may be distressed by the repetition.

The third process is a stable \emph{preference for familiar foods}, which is distinct from neophobia in that the child may be willing to taste something new but consistently prefers an established set. Such a preference is not pathological. The set of accepted foods in early childhood is narrow in most children and broadens with age, and an observer who measured only the breadth of the repertoire at a single time would conflate a child who is still broadening with a child whose repertoire has stalled.

The fourth process is \emph{variation in appetite}. Young children regulate their energy intake across days more accurately than they do within any single meal, so that a refused meal is often followed by a larger one. A behavior that appears as picky eating at one meal may be an appropriate response to satiety from a snack or a growth-related fluctuation in need. Treating each refusal as informative about preference, without attending to appetite, produces false positives.

The fifth process is \emph{resistance to coercion}. Evidence from several designs suggests that pressure to eat can run counter to its intended effect. Experimental work on restriction and pressure indicates that making a food conditional or forbidden increases its salience and desirability \citep{fisher1999restricting}. A longitudinal study of girls and their mothers found that mothers who ate more fruit and vegetables used less pressure and had daughters who were less picky and ate more fruit and vegetables \citep{galloway2005parental}, which associates pressure with pickiness without establishing whether pressure produces refusal, refusal elicits pressure, or both arise from another variable. Retrospective studies of forced consumption episodes find that adults frequently report lasting aversions to foods that they were compelled to eat as children \citep{batsell2002you}. Some portion of what presents as picky eating is therefore produced by the mealtime dynamics that adults introduce in response to earlier refusals, a feedback loop in which the remedy sustains the problem.

Beyond these five lies a sixth category that differs in kind: clinically significant feeding disorders. Avoidant/restrictive food intake disorder, as defined in the fifth edition of the American Psychiatric Association's diagnostic manual, describes persistent failure to meet nutritional needs associated with significant weight loss, nutritional deficiency, dependence on supplements, or psychosocial impairment, and is not explained by lack of food or by culturally sanctioned practice \citep{apa2013dsm5}. Children in this category need specialist assessment, and the techniques discussed in this essay are not a substitute for it. The disaggregation matters most here, because the same behavior that is a developmental stage in one child is a medical problem in another.

The disaggregation also yields a practical consequence that recurs in the sections that follow. Interventions are targeted at mechanisms, and a single outward behavior can have any of several mechanisms behind it. Repeated low-pressure exposure addresses neophobia. Textural accommodation addresses sensory aversion. Predictable meal structure and tolerance of variable appetite address the fourth process. Withdrawal of pressure addresses the fifth. A caregiver who applies the intervention for one mechanism to a child whose behavior arises from another will observe failure and may conclude that the child is intractable, when the more accurate conclusion is that the diagnosis of the refusal was wrong.

\section{What the Evidence Establishes: Exposure, Modeling, and Pressure}\label{sec:evidence}

Before introducing a new hypothesis it is necessary to state what the literature supports, since the hypothesis must be distinguished from it.

\subsection{Repeated exposure}

The most robust finding in the study of children's food acceptance is that repeated tasting of a novel food increases liking for and consumption of it. In an early experimental demonstration, \citet{birch1982dont} showed that two-year-old children's preferences shifted toward foods to which they had been exposed more often, an extension to food of the broader mere-exposure effect described by \citet{zajonc1968attitudinal}. Subsequent work has replicated the pattern across ages and foods, including with initially disliked vegetables in infants \citep{maier2007effects} and with a vegetable that children initially disliked \citep{wardle2003modifying}. \citet{cooke2007importance} reviewed laboratory and naturalistic studies and concluded that repeated opportunities to taste unfamiliar foods increase liking and consumption, and that familiarity is perhaps the most important determinant of a child's liking for a food. The number of exposures needed varies. In one study of infants, intake of an initially disliked vegetable rose to match that of a liked one by the eighth exposure \citep{maier2007effects}.

Two qualifications are important for the argument developed here. The first concerns the \emph{kind} of exposure. In a study of two- to five-year-olds that separated looking exposure from tasting exposure to novel fruits, \citet{birch1987what} found that looking raised visual preference judgments but did not raise taste preference judgments for foods that were never tasted, whereas tasting raised both, and concluded that experience must engage the modality relevant to the judgment. The result cautions against the assumption that familiarity with the appearance of a food converts into willingness to eat it. Later work complicates the picture. Picture books depicting foods increased toddlers' visual preferences for the pictured items, and for new pictures of them \citep{houstonprice2009picture}, and a nursery intervention in which children aged 12 to 36 months looked at, listened to, felt, and smelled unusual fruits and vegetables for four weeks increased their touching and tasting of those foods at a later meal, compared with a control group \citep{dazeley2015exposure}. Enjoyment of tactile play has also been associated with lower food neophobia in preschool children \citep{coulthard2015enjoyment}, though this association is correlational. The overall pattern is that non-taste exposure may help, that its effects are less consistent than those of tasting, and that the conditions under which it helps are not fully characterized.

The second qualification is that exposure is generally studied for a single food presented repeatedly in a fixed form. The literature has less to say about what happens when the form of the food is varied across exposures, or when acceptance of one food generalizes to related foods. There is some evidence of generalization. In a sample of infants aged four to seven months who were fed a target food daily for ten days, \citet{birch1998infants} found that intake of the target rose, that intake of the same food from another manufacturer and of similar foods also rose, and that intake of a different food did not change. This suggests that acceptance is partly organized by similarity and does not attach only to the specific item, though the sample was small, the infants were young, and the authors note that infants may have difficulty discriminating among many foods. This finding is the empirical foothold for the argument of Section~\ref{sec:hypothesis}, since it implies that the structure of similarity among foods matters to what children accept.

A third qualification concerns outcomes. Studies of exposure report increases in liking, in willingness to taste, and in intake, and these are related but not interchangeable. A child may consume a food more readily without rating it as more pleasant, and may rate it as more pleasant without eating more of it. Claims about the effectiveness of exposure should be read at the level at which the outcome was measured, and the same caution applies to the hypotheses developed below.

\subsection{Social modeling}

A second well-supported finding is that children's acceptance of foods is influenced by watching others eat. Early experimental work reported that an adult eating an unfamiliar food increased young children's willingness to taste it \citep{harper1975effect}. The manner of modeling matters. In a series of preschool quasi-experiments, silent teacher modeling did not encourage acceptance, whereas enthusiastic modeling maintained acceptance of a new food across five meals, and the addition of a competing peer model removed the effect \citep{hendy2000effectiveness}. A school-based program combining peer modeling, reward, and repeated exposure reduced neophobia in children \citep{laureati2014school}. The theoretical background is the social learning framework of \citet{bandura1977social}, according to which behavior is acquired in part through observation of models who are similar, competent, or valued.

The evidence for modeling is respectable but is not uniform, and several features limit its application. The effect of modeling depends on its manner, as the preschool studies above show, and in one of them acceptance of a new food was greatest at the first meal and then dropped. Moreover, modeling interacts with pressure. An adult who eats a food while simultaneously urging a child to eat it provides two signals, one of which may undermine the other. For present purposes the relevant point is that observation of a trusted person eating a food is a plausible route by which children come to regard the food as safe, and that this route does not require the child to eat anything at the moment of observation. The evidence on manner suggests that modeling is more effective when it is enthusiastic than when it is silent, which is a different thing from pressure.

\subsection{Pressure and restriction}

The third body of evidence, reviewed in Section~\ref{sec:disagg}, concerns the effects of coercion. The available evidence suggests that the direction of the effect is not merely neutral, although the designs differ in strength. Retrospective reports link forced consumption to lasting rejection of the food \citep{batsell2002you}, and longitudinal data link maternal pressure to picky eating in daughters \citep{galloway2005parental}. Neither design establishes causation, and the review of \citet{dovey2008food} lists pressure to eat among the factors that influence both neophobia and picky eating. Restricting access to appealing foods similarly increases the child's attention to and desire for them \citep{fisher1999restricting}. The practical lesson, which converges with the clinical responsive-feeding literature \citep{black2011responsive}, is that how a food is offered affects the response to it, independently of the food itself.

\subsection{What remains unexplained}

These three bodies of evidence together describe a set of effective conditions: repeated, pressure-free contact, accompanied by the example of a trusted eater. They leave open at least three questions that bear on the remainder of this essay. First, why does exposure sometimes fail, and what can be done in those cases other than to persist? Second, how does acceptance spread among foods, so that the work done on one food reduces the work required for another? Third, what role is played by the child's own control over the process, as opposed to the mere absence of pressure? The hypothesis that follows is an attempt to give a structured answer to the first two, and the discussion of agency in Section~\ref{sec:agency} addresses the third.

\section{Food Categories and the Structure of Familiarity}\label{sec:categories}

Sections~\ref{sec:evidence} to~\ref{sec:hypothesis} form the spine of the argument. The preceding section established what exposure and modeling can and cannot explain, this section supplies the account of food categories on which the hypothesis depends, and the next section states the hypothesis.

A food is never encountered as a bare collection of sensory properties. It is encountered as a member of a category, and the category determines what is expected of it. A child who has learned that crackers are crisp, savory, and dry has a set of expectations that is violated by a soft cracker, and the violation may be experienced as a reason for rejection even when the soft cracker would be perfectly acceptable if presented as something else. Understanding picky eating therefore requires a theory of how children structure food categories and how the boundaries of those categories are drawn.

\subsection{Prototypes and family resemblance}

The psychology of categorization offers two relevant ideas. The first is that natural categories are organized around prototypes, with membership graded by similarity to a central exemplar, and not defined by necessary and sufficient conditions \citep{rosch1975family}. The second, which precedes the psychological work, is the philosophical observation by \citet{wittgenstein1953philosophical} that many concepts are held together by overlapping similarities, a family resemblance, and not by a single shared feature. Applied to food, these ideas suggest that a child's category of, say, pizza is not a rule but a cluster. A prototypical pizza may have a round shape, a baked crust, tomato sauce, and melted cheese, and an item that shares several of these properties is a better or worse member of the category according to how many it shares.

If this is correct, then the boundaries of a food category are soft, and the child's willingness to accept a novel item depends on how well it fits into a category that is already accepted. This explains some familiar observations. A child who rejects a vegetable served alone may accept it as a topping on a food from an accepted category, and a child who rejects a homemade version of a familiar item may be responding to the failure of the homemade version to match the prototype, such as a differently shaped nugget or an unfamiliar brand of crackers.

\subsection{Food and non-food, and the semantics of disgust}

Categorization of food has a further dimension that makes it unlike the categorization of many other objects. Rozin and colleagues documented that the distinction between food and non-food is itself culturally constructed and affectively loaded, and that contamination and disgust attach to items that cross category boundaries \citep{rozin1980psychological,rozin1987perspective}. A food that looks like something inedible, or that contains an element perceived as contaminating, is rejected even when the child understands that it is safe. This implies that the regulation of food acceptance is not merely a matter of liking or disliking sensory properties but is also a matter of the child's judgment about what kind of thing a given item is.

The implication for transformation is double-edged. On the one hand, a preparation that resembles an accepted food on some dimensions may be treated as an instance of the accepted category and may inherit its acceptance. On the other hand, a preparation that violates category expectations on a dimension that matters to the child, for example a familiar dish with an unexpected color, may be treated as contaminated or wrong. Which dimensions matter is an empirical question, and it varies from child to child.

\subsection{Food chaining as a clinical precedent}

The idea that foods can be linked by shared properties and that acceptance can be extended along the links is not new to clinical practice. A technique called food chaining, described in a clinical guide by \citet{fraker2007food}, begins with foods that a child already accepts and proposes new foods that share sensory characteristics, such as a similar texture or flavor, moving gradually along dimensions of similarity. The technique is used by feeding therapists, and its rationale is similar to the one developed here. Its evidence base, however, is limited. A chart review of children treated with the approach reported favorable outcomes in children with feeding aversion \citep{fishbein2006food}, but a review of treated cases without a comparison group cannot separate the effect of the technique from maturation, repeated exposure, or the other components of treatment. No controlled trials are cited here, and the published guides do not provide a formal account of what counts as a small step or how dimensions should be weighted. The framework developed in the next section aims to supply the missing formal structure and to identify what an adequate test would look like, without claiming that the clinical technique has already been validated.

\section{The Transformational Familiarity Hypothesis}\label{sec:hypothesis}

\subsection{Preparations as vectors of properties}

To make the notion of similarity among foods precise, let a food preparation be represented as a vector of properties,
\begin{equation}
x = (f, t, a, v, c),
\end{equation}
where $f$ denotes flavor, $t$ texture, $a$ aroma, $v$ visual presentation, and $c$ contextual familiarity, meaning the setting, utensils, serving form, and social circumstances in which the preparation is encountered. Each component takes values in a space appropriate to it. Flavor, texture, and aroma may be described by categorical or ordinal descriptors, visual presentation by features such as shape and color, and context by features such as whether the food is eaten at the usual table and with the usual implements. No claim is made that these dimensions are exhaustive or independent. They are a working decomposition, chosen because caregivers commonly report that children react to each of them separately.

Let $\mathcal{X}$ denote the space of such preparations. At a given time, a child has an accepted set $S \subset \mathcal{X}$, consisting of preparations that the child will eat without hesitation, and an unaccepted region, consisting of preparations that the child refuses. Between them lies a boundary region of preparations that the child will explore under favorable conditions. The goal of an introduction is to expand $S$ to include a target preparation $z$ that the child currently refuses and that the adult wishes to make available, usually for nutritional or cultural reasons.

A dissimilarity measure on preparations is defined by
\begin{equation}
d_w(x, y) = \sum_{k \in \{f, t, a, v, c\}} w_k \, \delta_k(x_k, y_k),
\end{equation}
where $\delta_k$ is a dissimilarity on the $k$-th component and $w_k \geq 0$ are weights that reflect how much the child's response is governed by that dimension. The weights are child-specific. A child with strong textural sensitivity has a large $w_t$, and a child whose acceptance is largely governed by familiar settings has a large $w_c$. The weights are not assumed to be constant across foods or over time, and in the present formulation they are unknown quantities to be estimated from the child's behavior, not parameters supplied by the theory.

\subsection{Chains and the hypothesis}

\begin{definition}[Candidate chain]
A candidate chain from an accepted preparation $x_0 \in S$ to a target $z$ is a finite sequence $x_0, x_1, \ldots, x_n = z$ of distinct preparations such that (i) each consecutive pair is close under $d_w$, that is, $d_w(x_{i-1}, x_i) \leq \epsilon$ for some tolerance $\epsilon$; and (ii) each preparation is safe, nutritionally acceptable, and offered openly as a variant of the preceding one.
\end{definition}

Conditions (i) and (ii) are geometric and procedural. They do not guarantee that the child will accept any stage of the chain, still less its endpoint. A candidate chain is therefore a proposal about a path, and not a claim that the path works.

\begin{definition}[Monotone chain]
A candidate chain is monotone if the distance to the target never increases along it, that is, $d_w(x_i, z) \leq d_w(x_{i-1}, z)$ for every $i$.
\end{definition}

Monotonicity is a property of a restricted class of chains and not a requirement of the theory. It is a natural first model, since it excludes chains that wander, but it also excludes potentially useful detours. A child might need to become familiar with a new sauce on an otherwise familiar food before meeting the intended combination, and that intermediate preparation could lie farther from the target under the chosen distance while creating a more effective route toward acceptance. Geometric progress toward the target and psychological progress toward acceptance need not coincide, and whether monotone chains outperform nonmonotone ones is an empirical question that a later study could address.

To state what working means, the theory needs an object that depends on the child's history and not on the preparation alone.

Let $Y_j$ indicate whether the child voluntarily performs approach behavior $j$ (touching, smelling, tasting, swallowing, or requesting the food) toward a presented preparation in a given encounter, and let
\begin{equation}
A_j(x \mid h, s) = \Pr(Y_j = 1 \mid x, h, s)
\end{equation}
be the child's propensity to perform it, given the preparation $x$, the history $h$ of prior encounters, and the child's current state $s$, which includes appetite, mood, and relevant sensory conditions. The central object of the theory is a function of three arguments. Two identical preparations can elicit different responses because they were reached through different histories, and a theory that attended only to the preparation would miss exactly the effect at issue.

Because the propensity is latent, it cannot be observed directly, and an estimate of it from a few encounters is noisy. Success is therefore defined by an empirical criterion, and the propensity $A_j$ is reserved for the theoretical model.

\begin{definition}[Successful chain]
Fix in advance a criterion in the form of a number $k$ of voluntary performances of behavior $j$ within $m$ consecutive offers. A candidate chain is successful for behavior $j$ under this criterion if, when it is followed in order, the child meets the criterion at each intermediate stage before the next stage is introduced, and meets it at the final stage as well. The values of $k$ and $m$ are registered in advance, and no values are proposed here.
\end{definition}

Meeting the criterion is an observed event. It is evidence about $A_j$ and is not identical to it. Keeping the two apart avoids a circularity in which chains are called successful because an estimated propensity is high and are then used to confirm a theory about that propensity. Progress is defined by changes in observed behavior and not by geometric proximity alone. A candidate chain that stalls at stage $i$ has not been refuted as a geometric object. It has supplied information about the weights $w$, or about the dimension along which the step was too large.

The final clause of the candidate-chain definition is substantive. The chain is a sequence of transparent variations, and the child is not deceived about what is being offered. Concealing a target food inside an accepted one may produce immediate consumption, but it does not alter the child's categories, and it risks the loss of trust when the concealment is discovered. Whether concealment and transparency differ in their long-run effects on acceptance is itself an empirical question, though the clinical and ethical considerations favor transparency.

The hypothesis can now be stated.

\begin{hypothesis}[Transformational Familiarity]
Under conditions that are safe, voluntary, and free of pressure, and in which the child controls the pace and extent of participation, a target $z$, which is the same preparation in both cases, is more likely to be voluntarily approached and eventually accepted when it is reached by a candidate chain from an accepted preparation than when it is presented directly, holding constant the number of offers of the target food, the degree of adult modeling, time with the adult, and the child's baseline preferences. In the notation above, $A_{j^*}(z \mid h_{\text{chain}}, s) > A_{j^*}(z \mid h_{\text{direct}}, s)$ for the confirmatory behavior $j^*$ specified in the registered plan, with states $s$ matched.
\end{hypothesis}

No numerical value is asserted for either propensity, and none is available. The hypothesis is directional and comparative, and it stands or falls on the sign of the difference. It is a hypothesis about the effect of the history and not about the preparation, since $z$ is the same in both cases and only the path by which it was reached differs. The distinction between a candidate chain and a successful chain matters here, because the hypothesis concerns candidate chains in general. A study that evaluated only chains that happened to succeed would be selecting on the outcome.

Two refinements deserve separate statement.

\begin{hypothesis}[Participation]
Among children receiving the same chain, those who participate in the construction of the transformations are more likely to explore the target than those who receive the same preparations ready-made.
\end{hypothesis}

\begin{hypothesis}[Dimension specificity]\label{hyp:dim}
The tolerance $\epsilon$ that supports progress along a chain varies by dimension, and for some children and some dimensions, finer steps do not improve acceptance and may reduce it.
\end{hypothesis}

The third hypothesis guards against an assumption that is natural but unsupported, namely that smaller changes are always better. A child with strong sensory sensitivity may detect minute variations in texture, and a sequence of nearly identical but not quite identical preparations may be experienced as an unreliable series of near-misses and not as a gentle progression. In such a case the optimal strategy may be fewer, larger steps along some dimensions with strict constancy along others. The theory therefore does not specify a universal step size. It specifies a search problem in which the adult seeks, through observation of the child's responses, a path that holds fixed the properties the child cares most about while varying the properties the child tolerates varying.

A reasonable question is what the formalism adds to the ordinary-language statement that unfamiliar foods should be introduced gradually and openly. By itself, the notation is a way of stating a search problem, since the weights and the tolerance are not fixed by the theory. Its payoff is that it makes available claims the ordinary statement does not: that the dimensions of variation are separable, that children differ in their weights over them, and, as Hypothesis~\ref{hyp:dim} states, that the best step size is therefore specific to the dimension and to the child. The most direct test of the formalism is accordingly whether chains tailored to a child's estimated weights outperform chains with uniform steps, and the protocol of Section~\ref{sec:protocol} should be read with that contrast in mind.

A final hypothesis concerns a higher-order effect. A chain changes a child's encounters with one target, but it may also change the child's expectations about variation itself. A child who repeatedly meets familiar foods with small, openly explained changes might learn not only that a particular unfamiliar ingredient is acceptable but that modifications of familiar foods are generally safe to investigate.

\begin{hypothesis}[Transfer]\label{hyp:transfer}
Children who have followed a candidate chain toward a target are more likely than children who received matched direct presentation of that target to voluntarily explore a novel preparation unrelated to the training target, that is, $A_{j^*}(z' \mid h_{\text{chain}}, s) > A_{j^*}(z' \mid h_{\text{direct}}, s)$ for a transfer preparation $z'$ that was not part of the chain and does not share its key ingredients.
\end{hypothesis}

The first-order hypotheses concern acceptance of the particular target. The transfer hypothesis concerns something more general. If it held, it would suggest that the intervention changed the child's expectations about culinary variation, or the child's confidence in exploring it, and not only familiarity with particular ingredients. That would be a more distinctive contribution than any effect on a single target, since exposure to the target alone would not anticipate it.

\subsection{A worked example}

Consider a household in which a child reliably accepts crackers with cheese, and in which the adult wishes to extend the child's acceptance toward a pizza built on a crust of cooked pasta, a preparation that combines familiar elements (tomato sauce, cheese, baking) with a crust that is unfamiliar. The preparation is a pasta pizza that holds together as a pie. The chain proposed below proceeds through intermediate preparations, each of which differs from its predecessor in one or two properties. Table~\ref{tab:chain} describes the chain with qualitative descriptors. No numerical scores are assigned, since none would have an empirical basis.

\begin{table}[h]
\centering
\small
\renewcommand{\arraystretch}{1.25}
\begin{tabular}{@{}p{0.18\textwidth}p{0.14\textwidth}p{0.17\textwidth}p{0.16\textwidth}p{0.2\textwidth}@{}}
\toprule
Preparation & Base (texture) & Flavor & Appearance & Context \\
\midrule
$x_0$: cracker with cheese & crisp, dry & savory, mild cheese & small, flat, square & usual snack setting \\
$x_1$: cracker pizza & crisp, dry & adds tomato sauce & small, round, topped & usual snack setting \\
$x_2$: pita pizza & chewy, soft & same sauce and cheese & larger, round, topped & shared meal \\
$x_3$: English muffin pizza & soft, spongy & same sauce and cheese & thick, round, topped & shared meal \\
$x_4$: flatbread pizza & thin, baked & same sauce and cheese & large, cut in wedges & shared meal \\
$z$: pasta-crust pizza & softer, strands & same sauce and cheese & large, cut in wedges & shared meal \\
\bottomrule
\end{tabular}
\caption{An illustrative transformation chain. In every step the flavor core (tomato sauce and cheese) and the round, topped presentation are held constant while the base changes. The descriptors are qualitative, and the chain is an example of the form of the proposal, not a validated sequence.}
\label{tab:chain}
\end{table}

Several features of the example are instructive. In each transition from $x_{i-1}$ to $x_i$ the flavor core is preserved, so that the dimension $f$ is held constant while the base texture $t$ changes. At $x_1$ the change is the addition of a flavor that is familiar from other contexts. At $x_2$ the base becomes softer, which alters texture while keeping the surrounding properties intact. At $x_4$ the visual presentation changes through the larger size and the cut wedges, in preparation for the final step, at which the base changes from a baked bread to a pasta strand structure. The final preparation is identifiable as a variant of the first only because of the intermediate steps. Without them, the combination of pasta and pizza, two categories that the child may treat as mutually exclusive, would likely be encountered as a category violation.

The example also illustrates the sense in which the chain exploits family resemblance. No single property links the cracker with cheese to the pasta-crust pizza. The cracker pizza shares sauce and cheese with the pasta pizza, and it shares a crisp base with the original cracker, but the endpoints have almost nothing in common sensorially. The chain succeeds, if it does, because each adjacent pair overlaps enough that the child can treat the new preparation as another member of a familiar family.

The example also exposes the points at which the chain might fail. If the child is sensitive to the transition from crisp to chewy, then $x_2$ may be the stopping point, and no amount of repetition at later positions will compensate. If the child's acceptance of cracker pizza depends on the context of the snack, then moving to the shared meal at $x_2$ changes $c$ and $t$ simultaneously, and the effect cannot be attributed to either. The sensible response is to treat the chain as provisional and to adjust it according to the child's responses, using finer steps where the child resists and coarser steps where the child advances easily.

The chain in Table~\ref{tab:chain} is also contingent on a child profile. It assumes a child whose weight on texture is moderate. For a child with a high weight on texture, the transition from a crisp base to a chewy one at $x_2$ is a cliff and not a step, and a better chain would hold texture constant for as long as possible. Every base could be kept crisp, for example a cracker, a crisp flatbread, and a pasta crust baked until crisp, while the flavor and the context change, and the base could then be softened gradually, for instance by shortening the baking time, only at the end. The same target would then be reached by a different path, which is what Hypothesis~\ref{hyp:dim} predicts.

\subsection{Relation to exposure and generalization}

The hypothesis does not compete with the exposure literature, and it is important to specify how the two relate. Exposure effects predict that repeated tasting of $z$ will increase acceptance of $z$. The chain hypothesis predicts that, in addition, the \emph{route} by which $z$ is reached matters, and that an equal number of exposures to $z$ will be more effective when they are preceded by exposure to neighboring preparations that carry some of the properties of $z$. The generalization finding of \citet{birch1998infants} supports the plausibility of this, since it suggests that acceptance of one food transfers to similar foods. The chain hypothesis converts that observation into a design principle: it proposes constructing the similar foods deliberately and ordering them by their distance from what is already accepted.

The hypothesis also differs from simple exposure in the nature of the counterfactual. In a pure exposure design, the experimenter presents the same item repeatedly, and any resistance is overcome through accumulated familiarity. In the chain design, the item itself changes across presentations, and the aim is not to habituate the child to a fixed target but to bring the target within the family of accepted foods by changing the family boundary as well as the child's response. The two strategies can be combined, and a complete intervention would presumably include repeated exposure to the final preparation once it has been reached.

\subsection{Alternative explanations}\label{sec:alt}

An honest statement of the hypothesis must list the ways it could be wrong. The first alternative is that the chain works by exposure alone, since each step in the chain exposes the child to the target's component ingredients. If so, a design that matched exposure to those ingredients without the ordering would perform equally well. The second is that the chain works through social modeling, since adults who construct elaborate chains probably also eat more of the foods in question and attend more closely to the child. A third is that the benefit is due to the reduction of pressure, since an adult engaged in making variations is less likely to demand that the child eat. A fourth is a simple expectation effect: a parent who believes in the method will notice successes and ignore failures. Each of these is consistent with observed improvements and would not support the hypothesis. The experimental protocol of Section~\ref{sec:protocol} is organized to separate them.

\section{Observational and Participatory Familiarization}\label{sec:approach}

\subsection{Approach without ingestion}

A central assumption of the chain hypothesis is that the child's first successful interaction with a new preparation need not be consumption. Looking at, smelling, touching, helping to prepare, and tasting without swallowing are all forms of approach, and each can serve as a step toward eventual acceptance. This assumption is consistent with the clinical practice of feeding therapy, which commonly uses a graded sequence of interactions that begins with tolerating the presence of a food and ends with eating it. It is also consistent with the demand that the process be low-pressure, because a request to eat is a request for the most difficult step, and a child who is asked only to look at or touch a food is not placed in a position where refusal is the only available response.

The evidence on this point is, as noted in Section~\ref{sec:evidence}, mixed. The finding of \citet{birch1987what} that looking exposure did not raise taste preferences for foods that were never tasted stands against any simple claim that visual familiarity is sufficient. The nursery intervention of \citet{dazeley2015exposure}, by contrast, found that structured non-taste sensory exposure increased willingness to taste. A reasonable reconciliation is that passive looking does little, whereas active, multi-sensory interaction in a playful context does something, and that the benefit is mediated by the reduction of anxiety about the object and not by learning about its taste. The reconciliation is a conjecture, though it makes a testable prediction: the benefits of non-taste exposure should be larger for children with higher baseline neophobia and should be accompanied by changes in approach behavior that precede changes in tasting.

Approach behaviors can be defined for measurement, with the caveat that they need not form a single ladder. Candidates are tolerating the food's presence at the table, looking at it, touching it with a utensil or hand, smelling it, touching it to the lips, tasting without swallowing, swallowing a small amount, and eating a portion. Clinical practice often orders these from least to most demanding, but the ordering is an assumption. A child may smell a food and refuse to touch it, or taste it without wanting to handle it. Each behavior should therefore be recorded separately, as Section~\ref{sec:protocol} specifies, and the ordering treated as an empirical question. Separate recording still allows progress to be detected in children who have not yet eaten the food. This matters because outcomes confined to intake are insensitive to early change, and a study that measures only whether the child ate the food will classify as failures cases in which the child advanced in other ways.

\subsection{Modeling as a participatory structure}

The social modeling literature supports the use of a trusted adult who eats the target food in the child's presence. The way modeling is embedded in the interaction matters. When an adult eats a food as part of ordinary activity, the child receives information about the food's safety with no demand attached. The preschool experiments of \citet{hendy2000effectiveness} suggest that such modeling is more effective when it is enthusiastic than when it is silent, so demonstration should not be read as the absence of expression. When the adult eats while urging the child to do likewise, the information is accompanied by pressure, and the effect may be reversed. A relevant distinction is between modeling as demonstration, which may be warm and expressive, and modeling as persuasion, which is to say between communicating positive engagement and demanding compliance. The former is compatible with the responsive-feeding framework \citep{black2011responsive} and the latter is not.

Modeling has a further role in the chain framework. Because an adult who is making intermediate preparations is also eating them, the adult serves as a model for each step, and the child sees the same person eating a sequence of related foods, each of which is a variation of the last. This provides evidence, accessible to the child, that the variations belong to a single family and are safe within it. The point is again conjectural, but it is consistent with the finding that children are more willing to eat foods that they have seen a familiar person eat.

\section{Procedural and Causal Familiarity: Cultivation and Preparation}\label{sec:procedural}

\subsection{Three kinds of familiarity}

The discussion so far has treated familiarity as a matter of the sensory properties of a finished food. This is too narrow. A child may become familiar with a food in at least three distinct ways, and each makes a different contribution.

\emph{Sensory familiarity} arises from direct experience of a food's appearance, smell, texture, and taste. It is the kind of familiarity that the exposure literature measures, and it develops through repeated encounters with the food.

\emph{Procedural familiarity} arises from knowledge of how a food is obtained and prepared, including where it comes from, what is done to it, and what role the child's own actions play in the process. A child who has planted seeds, watered the plants, harvested the crop, washed and peeled the produce, and helped prepare it has a history with the food that extends well beyond the plate.

\emph{Causal familiarity} arises from understanding the relation between a raw material and a finished dish. A child who mixes flour, water, and yeast, observes the dough rise, and sees the loaf come from the oven learns that bread is the outcome of a process, and so does not meet it as an unexplained finished object. This kind of familiarity makes the transformations intelligible, and intelligibility is a precondition for a sense of control.

The three are not independent and need not co-occur. A child can be thoroughly familiar, sensorially, with a food whose origins are opaque, and a child can be procedurally familiar with a vegetable that has never been tasted. The point of the distinction is that each type of familiarity may reduce neophobia by a different mechanism, and that a food environment which supplies only the first has not exhausted the available routes.

\subsection{What the evidence shows}

Research on gardening and cooking interventions provides suggestive but limited support. A school-based garden curriculum was associated with improved knowledge of nutrition and preference for some vegetables in fourth-grade children \citep{morris2002garden}, and a pre-post study of middle school students reported improved recognition of, attitudes toward, preferences for, and willingness to taste vegetables \citep{ratcliffe2011effects}. A systematic review of fourteen gardening interventions reported that ten found statistically significant increases in children's fruit and vegetable consumption, described the evidence as mixed and limited by convenience samples, small samples, and self-reported intake, judged that most studies suggest a small positive influence, and called for controlled and randomized designs \citep{savoierosko2017increasing}. On the cooking side, a between-subject experiment with 47 children aged 6 to 10 found that children who helped a parent prepare a lunch ate more salad, more chicken, and more calories than children whose parent prepared it alone \citep{vanderhorst2014involving}, a result that is consistent with the participation hypothesis, though the sample was small, the children were older than most of those in the exposure literature, and the setting was a single meal.

The limitations of this evidence should be stated plainly. Participation in gardening and cooking is typically offered by motivated families and schools, so that children who participate may differ from those who do not. Interventions bundle several components, including nutrition education, enthusiasm of the adults, and the sheer fact of attention, and the studies do not isolate which component produces the effect. The question of whether participation changes acceptance of a food, as distinct from knowledge about it, is therefore open. What the evidence supports is a weaker claim. Several studies report favorable outcomes, but their designs do not establish the independent contribution of participation, and the absence of observed harm in them does not show that such programs are harmless in every setting. Effects on knowledge and attitudes are also more consistently reported than effects on intake.

\subsection{A proposed mechanism}

The three-way distinction suggests a mechanism by which participation could matter. A food that is encountered only as a finished commodity is, from the child's perspective, an opaque object that has arrived from somewhere unknown. Its safety is a matter of faith in the adult who serves it. A food whose production the child has watched, and in which the child has participated, is an object whose history is known. The child has seen it as a plant, as a raw ingredient, and as a result of specific actions, and has contributed to some of those actions. Neophobia, understood as caution toward substances whose properties are uncertain, may be reduced when the uncertainty is reduced by direct knowledge of the substance's origin and treatment.

On this account, participation reduces the novelty of the food in a way that mere tasting cannot, since it increases the number of channels through which the child has encountered the food and gives the child a role in producing it. It also increases the child's sense of authorship over the food, which may reduce the resistance that arises when food is imposed from outside. The mechanism would predict that participation matters more for foods that look unfamiliar in their finished state, such as whole vegetables with unusual shapes, than for foods that are already familiar in form.

This yields a refinement of the participation hypothesis. If participation works by increasing procedural and causal familiarity, then its effect should depend on the extent to which the child actually understands the process, and not merely on whether the child was present. A child who was assigned a trivial task and understood nothing about the larger operation should benefit less than a child who understood the sequence of transformations. The refinement is testable, since comprehension can be assessed independently of participation.

\subsection{Value beyond acceptance}

It would be a mistake to evaluate cultivation and cooking solely by their effects on what children agree to eat. Growing and preparing food teaches biological processes, measurement, sequencing, tool use, and practical competence, and it provides opportunities for sustained collaboration between adults and children. These benefits do not depend on the child's developing a taste for any particular vegetable. A program of participation that fails to alter acceptance of a target food may still be worthwhile on educational grounds, and a conclusion that treats acceptance as the only criterion of success would undervalue it. This observation is relevant to the framing of the argument as a whole. If food acceptance is treated as the only outcome that matters, the pressure to produce it returns by another route, and the conditions that support exploration are undermined.

\section{Structured Environments and the Distribution of Agency}\label{sec:agency}

\subsection{The division of responsibility}

The most influential practical framework for the relation between caregivers and children at mealtimes is the division of responsibility proposed by \citet{satter1986feeding}. On this account, the adult is responsible for what food is offered, when, and where, and the child is responsible for whether to eat and how much. The framework relies on the assumption that children, given a suitable environment, are capable of regulating their intake and of broadening their acceptance over time, and that pressure from the adult interferes with both. Its principles are consistent with the evidence on pressure and restriction reviewed in Section~\ref{sec:evidence}, and with the more general literature on responsive feeding, which treats feeding as an interaction in which the caregiver responds contingently to the child's signals \citep{black2011responsive}.

The framework is sometimes misread as permissiveness. It is not permissive about the adult's responsibilities, which include the planning, preparation, and presentation of food, the maintenance of a predictable schedule, and the modeling of eating. What it withholds from the adult is control over the child's individual decisions about eating. This is a different distribution of agency within a structured environment, not an absence of structure. The adult controls the conditions under which the child decides, and the child controls the decision.

\subsection{Autonomy and reactance}

Two bodies of psychological theory help to explain why this distribution is effective. The first is self-determination theory, which holds that sustained motivation depends in part on the satisfaction of a need for autonomy, the experience of acting from one's own volition \citep{deci2000what}. A child who eats because of external demands is not experiencing the behavior as self-directed, and the foods eaten under such conditions may acquire negative associations. The second is the theory of psychological reactance \citep{brehm1966theory}, according to which people respond to threats to their freedom by acting to restore it, often by doing the opposite of what is demanded. A demand to eat a food threatens the child's freedom to refuse, and refusal is the available means of restoring it. On this view, a fraction of what appears to be picky eating is reactance, and the appropriate intervention is the removal of the threat.

It would be a mistake, however, to treat every refusal as an exercise of autonomy. Some refusals arise from fear, some from sensory distress, and some from fatigue, and attributing them all to a desire for control is a form of romanticization. The claim of this section is narrower: whatever the underlying cause of a refusal, the child's retention of control over the decision removes one source of resistance, namely the resistance to being controlled, and leaves the other sources to be addressed by other means.

\subsection{Appetite as a facilitating condition}

A related question concerns the role of hunger. It is a common saying that hunger improves the taste of food, and there is a sound physiological basis for the idea that food has greater motivational value when one is ready to eat. Offering meals at times when children are likely to have an appetite is therefore sensible, and a snack that spoils appetite shortly before a meal works against exposure to novel foods.

The distinction that matters is between appetite as a facilitating condition and hunger as an instrument of coercion. The first consists in arranging the schedule so that the child arrives at meals with an appetite. The second consists in withholding food, or alternatives to a refused food, until the child's hunger is great enough that refusal becomes costly. These are different practices with different consequences. Deliberate deprivation, intended to make refusal untenable, is a form of pressure and is subject to the same adverse effects as other forms. A child who is very hungry is not in a good state for exploration, since distress is incompatible with the relaxed attention that approach requires. Young children generally do best with predictable opportunities to eat and the freedom to respond to their own hunger and fullness cues. The evidence on the direct effect of hunger on acceptance of novel foods is limited, and the recommendation to schedule meals sensibly rests on the broader responsive-feeding literature and not on a body of experiments specific to hunger.

The same caution applies to adults who manage their own eating by delaying meals until appetite is high. This may suit some individuals and some circumstances, but it is a personal pattern and not a general prescription, and its suitability depends on whether it supports adequate nutrition, energy, and well-being.

\section{Culinary Competence, Optionality, and the Economics of Preference}\label{sec:economics}

\subsection{The cost of flexibility}

The framework developed so far places demands on the adults who organize the meal. A household that is to offer intermediate preparations, vary the form of a food across meals, accommodate several preferences, and involve children in preparation needs to be able to cook, and to cook with some fluency. If every variation requires a new recipe, a new ingredient purchase, and a new set of instructions, the cost of flexibility is high, and a rigid menu of a few accepted dishes is the economical alternative. Households with limited time, equipment, or money may accordingly find it harder to maintain a varied repertoire, and it is plausible that this affects the range of foods their children are offered. Whether it also affects how often such children are described as picky is an empirical question on which this essay has no direct evidence.

This observation reframes part of the picky-eating problem as a problem of household production. Culinary competence changes the cost of variation, and changes it in a way that does not scale with the number of purchased ingredients. A cook who can turn an inexpensive ingredient into several preparations, adjust the seasoning of a basic product to suit different tastes, and combine leftovers into new dishes has a larger set of options from the same inventory than a cook who can only reheat.

\subsection{Culinary optionality}

The notion can be made precise. Let $I$ be the inventory of ingredients available to a household at a given time, let $\tau$ be the time available for preparation, let $K$ denote the cook's skills and equipment, and let $\mathcal{M}$ be the set of all meals. Because small changes in seasoning, portion size, or timing generate indefinitely many nominally different meals, counting meals directly would let optionality be inflated by trivial variation. Let $\sim_\rho$ be an equivalence relation on $\mathcal{M}$ under which two meals are equivalent if they fall within the same sensory and nutritional category at resolution $\rho$, and let $\mathcal{M}/{\sim_\rho}$ be the set of equivalence classes. Let $\mathcal{A}$ be the set of meals acceptable to the members of the household, taking account of preferences, dietary requirements, and the food-chaining state of each child. Define the \emph{culinary optionality} of the household at resolution $\rho$ as
\begin{equation}
O_\rho(I, \tau, K) = \bigl| \{ [m] \in \mathcal{M}/{\sim_\rho} : \text{some } m' \in [m] \text{ is producible from } I \text{ within } \tau \text{ given } K, \text{ and } m' \in \mathcal{A} \} \bigr|.
\end{equation}
The resolution $\rho$ is a modeling choice that must be stated with any measurement, and comparisons between households are meaningful only at the same resolution.
Optionality increases with the size of the inventory, but it also increases with skill and with the breadth of the acceptable set. Two households with identical inventories may differ greatly in optionality, because one can produce many acceptable meals from the ingredients and the other can produce only a few. Cooking competence expands the producibility condition without any change in expenditure.

Optionality has a direct application to picky eating. A household with high optionality can respond to a child's refusal by offering another preparation of the same nutritious ingredients, without needing to substitute different foods or to insist on the original. The response preserves the nutritional objective, preserves the child's opportunity to explore, and avoids the confrontation that arises when a single preparation is the only option. The limitation of a low-optionality household is not that its children are more difficult but that its structure leaves fewer ways to adapt to difficulty.

\subsection{The continuity of ingredients and the variability of meals}

A useful distinction is between repetition of ingredients and repetition of meals. A household can eat the same staple foods repeatedly without eating the same meal twice. Cooked rice may become fried rice on one day, a soup on another, and a stuffing on a third. A batch of beans may become a dip, a filling, a sauce, or a component of another dish. In each case there is material continuity, in that the same underlying ingredient is carried forward, and sensory discontinuity, in that the finished dish differs in flavor, texture, and appearance. The same noodles can participate in an Italian-style dish, a Southeast Asian-style stir-fry, and a curry without retaining the identity of the original.

This pattern bears on the theory in two ways. First, it shows that the identity of a dish is determined less by its ingredients taken individually than by the relations among them and the method of preparation, a point that parallels the family-resemblance account of categories. A child who rejects a vegetable in one dish may accept it in another, not because an aversion has been overcome but because the overall experience differs. Second, it provides a natural mechanism for chain construction. Yesterday's preparation is the starting material for today's, and the chain across days is partly a by-product of the way the household uses its inventory and not a program that must be separately designed.

Formally, let $I_t$ be the inventory at the start of day $t$, a vector of quantities of each ingredient and prepared item. During the day the household makes purchases $P_t$, selects a collection of items $X_t$ to be cooked, and converts them into prepared items $\Phi_t(X_t)$. The conversion map $\Phi_t$ records the change of identity that cooking effects, and it also absorbs losses such as trimming and evaporation, so that quantities are conserved only up to those losses. Let $E_t$ be the quantity eaten and $W_t$ the quantity discarded or otherwise lost. The inventory then evolves as
\begin{equation}
I_{t+1} = I_t + P_t - X_t + \Phi_t(X_t) - E_t - W_t,
\end{equation}
where each term is a vector indexed by item. Cooking therefore changes what the household possesses and not only how much of it: the inputs $X_t$ leave the inventory and the prepared items $\Phi_t(X_t)$ enter it. Prepared leftovers are entries of $I$ that are created in this way, and they carry their own identity, preparation state, and remaining safe storage life, which are tracked as attributes of the entry. A household that treats these entries as inputs to the next day's cooking, and not as waste, has a sequence of meals that is coupled through the inventory. The coupling reduces waste and cost, and it also produces a series of related preparations, which is the raw material of a transformation chain. The same representation lets culinary optionality be defined over a changing inventory, so that $O_\rho(I_t, \tau, K)$ varies from day to day as ingredients are transformed and consumed.

The storage-life attribute deserves separate emphasis, because it is what limits the length of a chain built from leftovers. Let $\sigma_i$ denote the remaining safe storage life of entry $i$ in the inventory. A step in a chain that draws on an earlier preparation must be taken while $\sigma_i$ remains positive, so a chain that depends on a single batch has a length bounded by the safe storage life of that batch divided by the interval between steps. Freezing portions at the outset extends the bound, at the cost of a change in texture that must itself be treated as a variation along the texture dimension. The constraint does not apply to chains whose steps are each cooked from fresh components, and the choice between the two is a trade-off among waste, cost, and flexibility. Public health guidance on safe storage periods governs the limits in practice, and no particular period is asserted here.

Optionality also enters the construction of chains directly. A candidate chain is a sequence of preparations that must each be producible from what the household has, so the number of candidate chains available from a given inventory grows with $O_\rho(I, \tau, K)$, and the cost of constructing a chain falls as optionality rises. The framework therefore predicts that chain-based interventions will be more feasible in households with higher optionality, and that optionality moderates the effect of the intervention. Both predictions can be tested by recording a household-level measure of optionality as a covariate in the protocol of Section~\ref{sec:protocol}.

The practice is subject to a constraint that must not be overlooked. Each transformation needs to preserve food safety, and repeated reheating does not reset the storage life of a cooked batch. Public health guidance generally advises that properly refrigerated cooked leftovers be used within a few days and that portions that will not be used in that time be frozen early. A strategy that relies on successive transformation of leftovers should be organized around these limits, and the household that follows it should freeze portions or discard food that has exceeded its safe storage period, not extend the sequence indefinitely.

\subsection{Improvisation and variety}

There is also a motivational dimension. Some cooks derive satisfaction from the process of solving a small, changing problem with whatever materials are at hand, and for them the changing inventory is a source of interest and not a burden. The preference for variety that this implies can be economically useful, since it favors flexible use of ingredients over fixed menus and so tends to reduce waste. A planning strategy in which the day's meal is determined jointly by the remaining shelf life of ingredients, the available time, and the preferences of the moment is the opposite of a rigid weekly plan, and its feasibility depends entirely on the cook's ability to improvise.

The household's decision problem can be expressed as the maximization, over a sequence of meals, of a weighted combination of enjoyment, variety, and nutritional adequacy, less waste and cost, subject to constraints on ingredients, time, equipment, and food safety. No claim is made that households actually solve such a problem or that the weights are stable. The formulation serves to display the structure of the trade-offs and to show that skill enters as a relaxation of constraints, and that variety and low cost are compatible given sufficient competence. It also shows why a household without that competence faces a harsher trade-off, in which variety can be bought only by spending more.

\subsection{Distributional limits}

These points must be qualified by recognition of the constraints that bound them. Time, equipment, food security, disability, and caregiver workload restrict how much experimentation a household can undertake. A caregiver working several jobs, or caring for several children, or lacking a functioning kitchen, cannot be expected to construct transformation chains. A strategy that presumes abundant time and skill, and then attributes the picky eating of children in resource-constrained households to a failure of method, would both mischaracterize the situation and add to the burdens that already exist. The argument of this section is that flexibility is valuable and that skill provides it at low cost. It is not that every household should be able to supply it. Any intervention based on the framework should therefore be designed with the constraints of the intended population in mind, and the research agenda should include simple versions of the strategies that require little time or equipment.

\section{Limits, Risks, and Misreadings}\label{sec:limits}

The argument of this essay is open to several misreadings that should be addressed directly.

The first is the inference that all food refusal is a solvable problem of presentation. It is not. Some children have persistent sensory, medical, or developmental feeding difficulties that require specialized assessment and support, and for these children a technique that works for ordinary neophobia may be ineffective or distressing. The framework offered here applies to the broad range of ordinary variation, and it comes with a recommendation that children who show weight loss, nutritional deficiency, severe restriction, distress at meals, or feeding difficulty that persists or worsens should be assessed by a clinician. It is also a reason for caution in using the chain approach with children whose sensory profile is unusual, since the preservation of fine sensory properties that such children require may leave little room for variation.

The second is the opposite inference, that every refusal expresses autonomy and must be respected as such. Section~\ref{sec:agency} argued that the retention of control by the child removes one source of resistance. It did not argue that the child's preferences are always correct or that adults have no role in setting the range of foods offered. The adult's responsibility for offering a nutritious and varied set of foods is retained.

The third misreading is that the transformation approach is a method of concealment. As specified in Section~\ref{sec:hypothesis}, the chain is a sequence of open variations, and the child is aware of what is offered. The distinction is both ethical and practical, since concealed foods do not alter the child's categories and may damage trust.

The fourth concerns safety. Allergies, choking hazards, and dietary requirements are not matters on which experimentation is appropriate, and the introduction of new foods must respect them. The framework applies to foods that are safe for the child, and the term voluntary exploration presupposes that the options offered are safe to explore.

The fifth concerns culture. The foods that a child is expected to accept, and the age at which acceptance is expected, differ across communities, and the notion of a picky eater is itself culturally situated. A framework that treats a narrow repertoire as a deficiency risks imposing the norms of one culinary tradition on another. Building chains from the household's own cuisine makes the construction of a chain culturally neutral in principle, but the choice of which foods to expand toward is not neutral, and it is made by the adults and the community, not by the framework. Which foods count as targets is therefore an input to the theory, and the argument has no implication that any particular foods should be preferred.

Scope is a further limitation. The essay concerns children, but category structure, procedural familiarity, and autonomy plausibly operate in adults as well, and food neophobia in adults is well documented \citep{pliner1992development}. The extension is a natural one, with the difference that adults can describe their own categories and consent to graded exposure, but it is not argued here.

Finally, there is the problem of attribution. Improvements in a child's acceptance of foods that follow the introduction of a new strategy are likely to be attributed to the strategy, though most children's repertoires broaden with age regardless. Without comparison conditions, the improvements cannot be assigned to the intervention, and the fact that the experience of caregivers is dominated by cases that did improve means that the method that was tried last will often receive the credit. This is the principal reason for insisting on the experimental design that follows.

\section{A Proposed Experimental Protocol}\label{sec:protocol}

The hypotheses of Section~\ref{sec:hypothesis} are testable, and this section outlines a design suited to the task. The design is a proposal, and its parameters would be refined in a pilot study before a full trial.

\subsection{Design}

The recommended design is a randomized comparison in a setting where meals are served regularly under a degree of experimental control, such as early-childhood education centers, with randomization at the level of the child or, where contamination between children is a concern, at the level of the classroom or center. Children are recruited with parental consent, with additional consent procedures appropriate to the setting, and children with diagnosed feeding disorders, relevant allergies, or medical dietary restrictions are excluded from the experimental conditions and referred to appropriate services. The protocol and analysis plan should be registered before data collection.

For each child, one or more target foods are selected from foods that the child currently refuses, as established in a baseline period by repeated offers without pressure. Selecting targets individually ensures that the test concerns foods that are actually unaccepted, and the use of several targets per child, analyzed with appropriate random effects, allows variation among foods to be modeled.

\subsection{Conditions}

The conditions are arranged so that the number and schedule of offers, adult modeling, time with the adult, and the child's motor and social engagement are matched as closely as possible across arms. Each child attends the same number of sessions of the same duration, the number of offers of each preparation is fixed in the protocol, and total adult attention is equalized by providing an attention-matched activity wherever a condition lacks a component. The matched activity is a non-food activity, such as handling modeling clay or tending a non-edible plant, supervised by the same adult with the same degree of social engagement. Because a fixed number of offers does not guarantee equal exposure, the number of tastes taken, the time the target food is within reach, and the amount of adult talk about the food are recorded and examined as mediators. Adult modeling is standardized in manner, with a script that is expressive but contains no invitation or pressure to eat, because the manner of modeling affects its effectiveness. Standardizing the script fixes its content but not its delivery, so the delivery is also measured. Sessions are recorded, and raters blind to condition code the adult's expressed enthusiasm, vocal and facial affect, eye contact with the child, and any invitation or pressure to eat. Enthusiasm may act through several routes, including attention, emotional arousal, perceived social approval, or the informational value of watching another person eat, and the coded measures allow the contribution of delivery to be examined as a moderator or mediator without assuming which route operates.

\begin{enumerate}[leftmargin=2em,itemsep=2pt]
\item \textbf{Direct presentation.} The target food is offered repeatedly in its standard form, with adult modeling and no pressure, together with the attention-matched activity. This is the standard exposure condition.
\item \textbf{Ordered chain, passive.} The child receives a sequence of ready-made preparations along a candidate chain from an accepted food to the target, in order, with the same number of offers and the same adult modeling as in condition 1. The attention-matched activity fills the time that preparation would otherwise occupy.
\item \textbf{Randomized-order chain, passive.} The child receives the same set of intermediate preparations as in condition 2, with the same number of offers, but in randomized order.
\item \textbf{Ordered chain, participatory.} The child helps to make the same sequence of preparations as in condition 2, so that the finished foods are equivalent and only the experience of construction differs. Children in this condition also receive the matched non-food plant activity.
\item \textbf{Ordered chain, participatory, with cultivation.} As in condition 4, with the addition that the child participates in growing or handling the raw ingredient of the target food, independently of the preparations, in place of the matched plant activity.
\item \textbf{Ingredient-matched control.} The child receives the same number of offers of the target's component ingredients in unrelated preparations that do not form a chain.
\end{enumerate}

Comparison of conditions 1 and 2 tests the central hypothesis, with attention and schedule matched. Comparison of 2 and 3 tests whether the order of the intermediate preparations matters, apart from exposure to the same set of preparations. Comparison of 2 and 6 tests whether the similarity structure of the intermediate preparations matters, apart from exposure to the target's ingredients, which addresses the first of the alternative explanations of Section~\ref{sec:alt}. Comparison of 2 and 4 tests the participation hypothesis, with the passive group receiving the matched activity so that motor activity, social engagement, and adult attention are equalized across the two arms. Comparison of 4 and 5 tests the incremental contribution of cultivation. A six-arm design is demanding, and the contrasts differ in priority. The comparisons of conditions 1 with 2 and of 2 with 6 are essential, since they test the central hypothesis and separate it from exposure to the ingredients alone, and the comparison of 2 with 3 is nearly essential, since it tests whether order matters. The participation contrasts (2 with 4, and 4 with 5) address a separate hypothesis and could be run as a second study, once the first has established whether chains work at all. A further contrast worth reserving for later is between chains tailored to a child's estimated weights and chains with uniform steps, which is the most direct test of the dimension-specificity hypothesis. Matching reduces confounding but does not eliminate it, since an attention-matched activity cannot be guaranteed to draw the same attention as the activity it replaces. The design also does not by itself separate the effect of modeling from the effect of reduced pressure, and for that purpose the adult's behavior should be recorded and rated by observers blind to condition wherever possible.

\subsection{Outcomes}

Outcomes are recorded separately and are not collapsed into a single scale, because the behaviors that mark approach need not form a universal hierarchy. A child may smell a food and refuse to touch it, or taste it without wanting to handle it. For each offer, the following are recorded as separate binary indicators or frequencies: touching, smelling, mouth contact or tasting, swallowing, and spontaneous requests for or questions about the food. The confirmatory outcome is voluntary tasting per offer, because it is the behavior most directly tied to the hypotheses and is unambiguous to code. The remaining behaviors are secondary outcomes, each analyzed with mixed-effects models for binary or count data, with random effects for child, food, and site, and with adjustment for multiplicity specified in the registered plan. A composite ordinal approach scale is retained as an exploratory outcome only, and its use is conditional on evidence from scaling analyses of the pilot data that the behaviors are in fact ordered in a consistent way. Three further outcomes are kept distinct because they are related but not interchangeable: liking, measured with an age-appropriate hedonic rating; willingness to taste, measured by behavior when the food is offered; and intake, measured in a standard serving. Parent-reported acceptance at home is a further secondary outcome. A follow-up measurement, taken after the intervention has ended, is needed to assess durability, since a gain that disappears when the experimental conditions are withdrawn would be of limited practical value. A chain is regarded as having reached its endpoint when the child shows stable voluntary exploration of the target across a number of consecutive offers specified in the registered plan, and acceptance in the sense of regular eating is analyzed as a separate, later outcome.

A transfer task tests Hypothesis~\ref{hyp:transfer}. After the intervention, each child is offered a novel preparation that is unrelated to the training target and does not share its key ingredients, presented openly as a variation of a different accepted food. The same behaviors are recorded, and voluntary exploration of the transfer preparation is compared across conditions. A difference between the chain conditions and the direct and ingredient-matched conditions would support the transfer hypothesis. An absence of difference would leave any first-order effect on the training target as the only candidate result.

\subsection{Covariates and moderators}

Baseline neophobia, measured with established child measures \citep{pliner1994development}, age, sensory sensitivity, the child's baseline repertoire, and a household-level measure of culinary optionality should be recorded. The theory predicts that effects will be moderated by these variables. In particular, dimension specificity (Hypothesis~\ref{hyp:dim}) predicts that the optimal chain differs by sensory profile, and the participation refinement of Section~\ref{sec:procedural} predicts that comprehension of the process moderates the effect of participation. Comprehension can be assessed by asking children to describe or sequence the steps they have performed.

\subsection{Analysis and interpretation}

The sample size should be determined by a power analysis based on pilot estimates of between-condition differences and variance components, and no sample size is asserted here. Analysis should follow the registered plan, with intention-to-treat as the primary approach. The interpretation should be proportionate to the results. A result in which condition 2 exceeds condition 1 but does not exceed condition 6 would be consistent with an explanation in terms of exposure to the ingredients, without establishing it, and the comparison of conditions 2 and 3 is the cleaner test of whether the order of the transformations matters. A result in which no condition exceeds the standard exposure control would count against the hypothesis. The hypothesis is intended to be refutable, and a protocol that could not produce a negative result would not be a test.

\subsection{Ethical considerations}

The study involves young children and food, and the standard protections apply, including parental consent, the child's ongoing assent as manifested in willingness to participate, the absence of pressure, and the freedom to decline any step without consequence. Foods are screened for allergens and choking risk. Because the protocol rests on the claim that coercion is counterproductive, any element of the design that created pressure to eat would contradict its premises, and staff should be trained accordingly.

\section{Conclusion}\label{sec:conclusion}

The argument of this essay has been that the question of how to make a child eat is badly posed. It treats the child as the site of the problem and the food as fixed, whereas the evidence points to a system in which the form of the food, the way it is introduced, the state of the child, and the structure of the household all contribute. A refusal conveys information about the interaction among these elements, and its interpretation depends on which of several mechanisms is operating. Where the mechanism is neophobia, repeated low-pressure contact is supported by a substantial literature. Where the mechanism is sensory aversion, accommodation is required. Where the mechanism is resistance to coercion, the remedy is the removal of pressure. Where the mechanism is a clinical feeding disorder, specialist care is needed.


The contribution can be stated compactly. The repeated-exposure model changes the eater's response to a fixed food. The proposal developed here changes the pathway by which the eater encounters the food, and treats the history of encounters, and not the food alone, as what determines the response. The first is well supported. The second is a conjecture that the protocol above is designed to test. A higher-order version of the conjecture, stated as the transfer hypothesis, holds that experience with chains may change expectations about variation itself, and it is tested by a transfer task. A related refinement concerns the endpoint of a chain. A chain need not terminate in consumption. It may terminate in a stable willingness to explore the target, with acceptance in the sense of regular eating treated as a separate and later outcome. Defining success in this way is consistent with the commitment to the child's control over participation, and it avoids classifying every attempt that stops short of eating as a failed intervention.
Beyond the established findings, the essay proposed that acceptance can be extended by building intermediate preparations that connect accepted and unaccepted foods, and that participation in the production of these preparations, and in the cultivation of their ingredients, may contribute procedural and causal familiarity in addition to the sensory familiarity that exposure supplies. This is the Transformational Familiarity Hypothesis, together with its participation and dimension-specificity refinements. These are conjectures. They draw plausibility from the generalization of acceptance among similar foods, from the structure of food categories, and from the clinical practice of food chaining, and they are untested in the form stated here. The protocol of Section~\ref{sec:protocol} specifies how they could be tested and what results would count against them.

The argument also relocates part of the problem. If food acceptance depends on the flexibility of the surrounding environment as well as on the flexibility of the eater, then households with greater culinary optionality, achieved through skill and not through expenditure, are better placed to support exploration, and households that lack the time or resources to develop it face a structural disadvantage that is easily mistaken for a feature of their children. The practical implication is that support for families should include the means of cooking flexibly, and that the evaluation of any strategy should consider what it requires of caregivers.

There is a last point that concerns the aim of the exercise. The value of involving children in the growing and making of food cannot be measured entirely by what they afterward agree to eat. A child who learns how a vegetable grows, how dough rises, and how a leftover can become a new dish has acquired understanding and competence that are worth having whether or not the vegetable becomes a favorite. If acceptance follows, it is a benefit. If it does not, the child has still gained something, and the conditions that make exploration possible, which are safety, control, and intelligibility, have been preserved. A framework that keeps those conditions in view is more likely to produce children who are willing to explore than one that treats each meal as a contest.

\begin{thebibliography}{99}
\setlength{\itemsep}{0.2em}

\bibitem[{American Psychiatric Association(2013)}]{apa2013dsm5}
American Psychiatric Association. (2013).
\newblock \emph{Diagnostic and Statistical Manual of Mental Disorders} (5th ed.).
\newblock Arlington, VA: American Psychiatric Publishing.

\bibitem[Bandura(1977)]{bandura1977social}
Bandura, A. (1977).
\newblock \emph{Social Learning Theory}.
\newblock Englewood Cliffs, NJ: Prentice Hall.

\bibitem[Batsell et~al.(2002)]{batsell2002you}
Batsell, W.~R., Brown, A.~S., Ansfield, M.~E., and Paschall, G.~Y. (2002).
\newblock ``You will eat all of that!'': A retrospective analysis of forced consumption episodes.
\newblock \emph{Appetite}, 38(3), 211--219.

\bibitem[Birch and Marlin(1982)]{birch1982dont}
Birch, L.~L. and Marlin, D.~W. (1982).
\newblock I don't like it; I never tried it: Effects of exposure on two-year-old children's food preferences.
\newblock \emph{Appetite}, 3(4), 353--360.

\bibitem[Birch et~al.(1987)]{birch1987what}
Birch, L.~L., McPhee, L., Shoba, B.~C., Pirok, E., and Steinberg, L. (1987).
\newblock What kind of exposure reduces children's food neophobia? Looking vs. tasting.
\newblock \emph{Appetite}, 9(3), 171--178.

\bibitem[Birch et~al.(1998)]{birch1998infants}
Birch, L.~L., Gunder, L., Grimm-Thomas, K., and Laing, D.~G. (1998).
\newblock Infants' consumption of a new food enhances acceptance of similar foods.
\newblock \emph{Appetite}, 30(3), 283--295.

\bibitem[Black and Aboud(2011)]{black2011responsive}
Black, M.~M. and Aboud, F.~E. (2011).
\newblock Responsive feeding is embedded in a theoretical framework of responsive parenting.
\newblock \emph{Journal of Nutrition}, 141(3), 490--494.

\bibitem[Brehm(1966)]{brehm1966theory}
Brehm, J.~W. (1966).
\newblock \emph{A Theory of Psychological Reactance}.
\newblock New York: Academic Press.

\bibitem[Cermak et~al.(2010)]{cermak2010food}
Cermak, S.~A., Curtin, C., and Bandini, L.~G. (2010).
\newblock Food selectivity and sensory sensitivity in children with autism spectrum disorders.
\newblock \emph{Journal of the American Dietetic Association}, 110(2), 238--246.

\bibitem[Cooke(2007)]{cooke2007importance}
Cooke, L. (2007).
\newblock The importance of exposure for healthy eating in childhood: A review.
\newblock \emph{Journal of Human Nutrition and Dietetics}, 20(4), 294--301.

\bibitem[Cooke et~al.(2007)]{cooke2007genetic}
Cooke, L.~J., Haworth, C.~M.~A., and Wardle, J. (2007).
\newblock Genetic and environmental influences on children's food neophobia.
\newblock \emph{American Journal of Clinical Nutrition}, 86(2), 428--433.

\bibitem[Coulthard and Thakker(2015)]{coulthard2015enjoyment}
Coulthard, H. and Thakker, D. (2015).
\newblock Enjoyment of tactile play is associated with lower food neophobia in preschool children.
\newblock \emph{Journal of the Academy of Nutrition and Dietetics}, 115(7), 1134--1140.

\bibitem[Dazeley and Houston-Price(2015)]{dazeley2015exposure}
Dazeley, P. and Houston-Price, C. (2015).
\newblock Exposure to foods' non-taste sensory properties: A nursery intervention to increase children's willingness to try fruit and vegetables.
\newblock \emph{Appetite}, 84, 1--6.

\bibitem[Deci and Ryan(2000)]{deci2000what}
Deci, E.~L. and Ryan, R.~M. (2000).
\newblock The ``what'' and ``why'' of goal pursuits: Human needs and the self-determination of behavior.
\newblock \emph{Psychological Inquiry}, 11(4), 227--268.

\bibitem[Dovey et~al.(2008)]{dovey2008food}
Dovey, T.~M., Staples, P.~A., Gibson, E.~L., and Halford, J.~C.~G. (2008).
\newblock Food neophobia and `picky/fussy' eating in children: A review.
\newblock \emph{Appetite}, 50(2--3), 181--193.

\bibitem[Fishbein et~al.(2006)]{fishbein2006food}
Fishbein, M., Cox, S., Swenny, C., Mogren, C., Walbert, L., and Fraker, C. (2006).
\newblock Food chaining: A systematic approach for the treatment of children with feeding aversion.
\newblock \emph{Nutrition in Clinical Practice}, 21(2), 182--184.

\bibitem[Fisher and Birch(1999)]{fisher1999restricting}
Fisher, J.~O. and Birch, L.~L. (1999).
\newblock Restricting access to palatable foods affects children's behavioral response, food selection, and intake.
\newblock \emph{American Journal of Clinical Nutrition}, 69(6), 1264--1272.

\bibitem[Fraker et~al.(2007)]{fraker2007food}
Fraker, C., Fishbein, M., Cox, S., and Walbert, L. (2007).
\newblock \emph{Food Chaining: The Proven 6-Step Plan to Stop Picky Eating, Solve Feeding Problems, and Expand Your Child's Diet}.
\newblock New York: Marlowe \& Company.

\bibitem[Galloway et~al.(2005)]{galloway2005parental}
Galloway, A.~T., Fiorito, L., Lee, Y., and Birch, L.~L. (2005).
\newblock Parental pressure, dietary patterns, and weight status among girls who are ``picky eaters''.
\newblock \emph{Journal of the American Dietetic Association}, 105(4), 541--548.

\bibitem[Harper and Sanders(1975)]{harper1975effect}
Harper, L.~V. and Sanders, K.~M. (1975).
\newblock The effect of adults' eating on young children's acceptance of unfamiliar foods.
\newblock \emph{Journal of Experimental Child Psychology}, 20(2), 206--214.

\bibitem[Hendy and Raudenbush(2000)]{hendy2000effectiveness}
Hendy, H.~M. and Raudenbush, B. (2000).
\newblock Effectiveness of teacher modeling to encourage food acceptance in preschool children.
\newblock \emph{Appetite}, 34(1), 61--76.

\bibitem[Houston-Price et~al.(2009)]{houstonprice2009picture}
Houston-Price, C., Burton, E., Hickinson, R., Inett, J., Moore, E., Salmon, K., and Shiba, P. (2009).
\newblock Picture book exposure elicits positive visual preferences in toddlers.
\newblock \emph{Journal of Experimental Child Psychology}, 104(1), 89--104.

\bibitem[Laureati et~al.(2014)]{laureati2014school}
Laureati, M., Bergamaschi, V., and Pagliarini, E. (2014).
\newblock School-based intervention with children: Peer-modeling, reward and repeated exposure reduce food neophobia.
\newblock \emph{Appetite}, 83, 26--32.

\bibitem[Maier et~al.(2007)]{maier2007effects}
Maier, A., Chabanet, C., Schaal, B., Issanchou, S., and Leathwood, P. (2007).
\newblock Effects of repeated exposure on acceptance of initially disliked vegetables in 7-month old infants.
\newblock \emph{Food Quality and Preference}, 18(8), 1023--1032.

\bibitem[Morris and Zidenberg-Cherr(2002)]{morris2002garden}
Morris, J.~L. and Zidenberg-Cherr, S. (2002).
\newblock Garden-enhanced nutrition curriculum improves fourth-grade school children's knowledge of nutrition and preference for some vegetables.
\newblock \emph{Journal of the American Dietetic Association}, 102(1), 91--93.

\bibitem[Pliner(1994)]{pliner1994development}
Pliner, P. (1994).
\newblock Development of measures of food neophobia in children.
\newblock \emph{Appetite}, 23(2), 147--163.

\bibitem[Pliner and Hobden(1992)]{pliner1992development}
Pliner, P. and Hobden, K. (1992).
\newblock Development of a scale to measure the trait of food neophobia in humans.
\newblock \emph{Appetite}, 19(2), 105--120.

\bibitem[Ratcliffe et~al.(2011)]{ratcliffe2011effects}
Ratcliffe, M.~M., Merrigan, K.~A., Rogers, B.~L., and Goldberg, J.~P. (2011).
\newblock The effects of school garden experiences on middle school-aged students' knowledge, attitudes, and behaviors associated with vegetable consumption.
\newblock \emph{Health Promotion Practice}, 12(1), 36--43.

\bibitem[Rosch and Mervis(1975)]{rosch1975family}
Rosch, E. and Mervis, C.~B. (1975).
\newblock Family resemblances: Studies in the internal structure of categories.
\newblock \emph{Cognitive Psychology}, 7(4), 573--605.

\bibitem[Rozin(1976)]{rozin1976selection}
Rozin, P. (1976).
\newblock The selection of foods by rats, humans, and other animals.
\newblock In J.~S. Rosenblatt, R.~A. Hinde, E.~Shaw, and C.~Beer (Eds.), \emph{Advances in the Study of Behavior} (Vol.~6, pp.~21--76). New York: Academic Press.

\bibitem[Rozin and Fallon(1980)]{rozin1980psychological}
Rozin, P. and Fallon, A.~E. (1980).
\newblock The psychological categorization of foods and non-foods: A preliminary taxonomy of food rejections.
\newblock \emph{Appetite}, 1(3), 193--201.

\bibitem[Rozin and Fallon(1987)]{rozin1987perspective}
Rozin, P. and Fallon, A.~E. (1987).
\newblock A perspective on disgust.
\newblock \emph{Psychological Review}, 94(1), 23--41.

\bibitem[Satter(1986)]{satter1986feeding}
Satter, E.~M. (1986).
\newblock The feeding relationship.
\newblock \emph{Journal of the American Dietetic Association}, 86(3), 352--356.

\bibitem[Savoie-Roskos et~al.(2017)]{savoierosko2017increasing}
Savoie-Roskos, M.~R., Wengreen, H., and Durward, C. (2017).
\newblock Increasing fruit and vegetable intake among children and youth through gardening-based interventions: A systematic review.
\newblock \emph{Journal of the Academy of Nutrition and Dietetics}, 117(2), 240--250.

\bibitem[van~der Horst et~al.(2014)]{vanderhorst2014involving}
van~der Horst, K., Ferrage, A., and Rytz, A. (2014).
\newblock Involving children in meal preparation: Effects on food intake.
\newblock \emph{Appetite}, 79, 18--24.

\bibitem[Wardle et~al.(2003)]{wardle2003modifying}
Wardle, J., Cooke, L.~J., Gibson, E.~L., Sapochnik, M., Sheiham, A., and Lawson, M. (2003).
\newblock Increasing children's acceptance of vegetables: A randomized trial of parent-led exposure.
\newblock \emph{Appetite}, 40(2), 155--162.

\bibitem[Wittgenstein(1953)]{wittgenstein1953philosophical}
Wittgenstein, L. (1953).
\newblock \emph{Philosophical Investigations}.
\newblock Oxford: Basil Blackwell.

\bibitem[Zajonc(1968)]{zajonc1968attitudinal}
Zajonc, R.~B. (1968).
\newblock Attitudinal effects of mere exposure.
\newblock \emph{Journal of Personality and Social Psychology}, 9(2, Pt.~2), 1--27.

\end{thebibliography}

\end{document}
